
\documentclass{article}
\usepackage{colortbl}
\usepackage{makecell}
\usepackage{multirow}
\usepackage{supertabular}

\begin{document}

\newcounter{utterance}

\centering \large Interaction Transcript for game `cladder', experiment `full\_v1.5\_default', episode 1133 with LLP: Large Language Problems\_\_qwen3.5{-}9b\_\_llp{-}final.
\vspace{24pt}

{ \footnotesize  \setcounter{utterance}{1}
\setlength{\tabcolsep}{0pt}
\begin{supertabular}{c@{$\;$}|p{.15\linewidth}@{}p{.15\linewidth}p{.15\linewidth}p{.15\linewidth}p{.15\linewidth}p{.15\linewidth}}
    \# & \multicolumn{2}{c}{Player} && \multicolumn{2}{c}{Game Master} \\
    \hline

    \theutterance \stepcounter{utterance}  
    & & & \multicolumn{4}{p{0.6\linewidth}}{
        \cellcolor[rgb]{0.9,0.9,0.9}{
            \makecell[{{p{\linewidth}}}]{
                \texttt{\tiny{[P1$\langle$GM]}}
                \texttt{You are an expert in causal inference. The following question is not a typical commonsense query, but rather a meticulously designed question created by a professor specializing in causal inference, intended to assess the students' mastery of the course content.} \\
\texttt{Imagine a self{-}contained, hypothetical world with only the following conditions, and without any unmentioned factors or causal relationships: Kidney stone size has a direct effect on treatment and recovery. Treatment has a direct effect on recovery. For patients who have small kidney stones and not receiving treatment, the probability of recovery is 3\%. For patients who have small kidney stones and receiving treatment, the probability of recovery is 33\%. For patients who have large kidney stones and not receiving treatment, the probability of recovery is 60\%. For patients who have large kidney stones and receiving treatment, the probability of recovery is 93\%. The overall probability of large kidney stone is 49\%. Will receives treatment increase the chance of recovery?} \\
\texttt{Start your answer with 'Yes' or 'No', followed by additional reasoning or evidence to support your explanation.} \\
            }
        }
    }
     \\ \\

    \theutterance \stepcounter{utterance}  
    & \multicolumn{4}{p{0.6\linewidth}}{
        \cellcolor[rgb]{0.561,0.737,0.561}{
            \makecell[{{p{\linewidth}}}]{
                \texttt{\tiny{[P1$\rangle$GM]}}
                \texttt{No.} \\
\\ 
\texttt{To determine whether treatment increases the chance of recovery for Will, we need to compare the probability of recovery under treatment versus no treatment, accounting for the distribution of kidney stone sizes.} \\
\\ 
\texttt{Let’s denote:} \\
\texttt{{-} \textbackslash{}( P(R|T, S) \textbackslash{}) = probability of recovery given treatment and stone size S} \\
\texttt{{-} \textbackslash{}( P(R|\textbackslash{}neg T, S) \textbackslash{}) = probability of recovery given no treatment and stone size S} \\
\texttt{{-} \textbackslash{}( P(S) \textbackslash{}) = probability of stone size S} \\
\\ 
\texttt{Given:} \\
\texttt{{-} \textbackslash{}( P(R|T, \textbackslash{}text\{small\}) = 0.33 \textbackslash{})} \\
\texttt{{-} \textbackslash{}( P(R|\textbackslash{}neg T, \textbackslash{}text\{small\}) = 0.03 \textbackslash{})} \\
\texttt{{-} \textbackslash{}( P(R|T, \textbackslash{}text\{large\}) = 0.93 \textbackslash{})} \\
\texttt{{-} \textbackslash{}( P(R|\textbackslash{}neg T, \textbackslash{}text\{large\}) = 0.60 \textbackslash{})} \\
\texttt{{-} \textbackslash{}( P(\textbackslash{}text\{large\}) = 0.49 \textbackslash{}), so \textbackslash{}( P(\textbackslash{}text\{small\}) = 0.51 \textbackslash{})} \\
\\ 
\texttt{The overall probability of recovery with treatment:} \\
\texttt{\textbackslash{}{[}} \\
\texttt{P(R|T) = P(R|T, \textbackslash{}text\{small\}) \textbackslash{}cdot P(\textbackslash{}text\{small\}) + P(R|T, \textbackslash{}text\{large\}) \textbackslash{}cdot P(\textbackslash{}text\{large\})} \\
\texttt{\textbackslash{}{]}} \\
\texttt{\textbackslash{}{[}} \\
\texttt{= 0.33 \textbackslash{}cdot 0.51 + 0.93 \textbackslash{}cdot 0.49 = 0.1683 + 0.4557 = 0.624} \\
\texttt{\textbackslash{}{]}} \\
\\ 
\texttt{The overall probability of recovery without treatment:} \\
\texttt{\textbackslash{}{[}} \\
\texttt{P(R|\textbackslash{}neg T) = P(R|\textbackslash{}neg T, \textbackslash{}text\{small\}) \textbackslash{}cdot P(\textbackslash{}text\{small\}) + P(R|\textbackslash{}neg T, \textbackslash{}text\{large\}) \textbackslash{}cdot P(\textbackslash{}text\{large\})} \\
\texttt{\textbackslash{}{]}} \\
\texttt{\textbackslash{}{[}} \\
\texttt{= 0.03 \textbackslash{}cdot 0.51 + 0.60 \textbackslash{}cdot 0.49 = 0.0153 + 0.294 = 0.3093} \\
\texttt{\textbackslash{}{]}} \\
\\ 
\texttt{Since \textbackslash{}( P(R|T) = 0.624 > P(R|\textbackslash{}neg T) = 0.3093 \textbackslash{}), treatment increases the overall chance of recovery.} \\
\\ 
\texttt{Thus, yes, receiving treatment increases the chance of recovery for Will.} \\
            }
        }
    }
    & & \\ \\

    \theutterance \stepcounter{utterance}  
    & & & \multicolumn{2}{p{0.3\linewidth}}{
        \cellcolor[rgb]{0.9,0.9,0.9}{
            \makecell[{{p{\linewidth}}}]{
                \texttt{\tiny{[GM$|$GM]}}
                \texttt{No} \\
            }
        }
    }
    & & \\ \\

    \theutterance \stepcounter{utterance}  
    & & & \multicolumn{2}{p{0.3\linewidth}}{
        \cellcolor[rgb]{0.9,0.9,0.9}{
            \makecell[{{p{\linewidth}}}]{
                \texttt{\tiny{[GM$|$GM]}}
                \texttt{yes} \\
            }
        }
    }
    & & \\ \\

    \theutterance \stepcounter{utterance}  
    & & & \multicolumn{2}{p{0.3\linewidth}}{
        \cellcolor[rgb]{0.9,0.9,0.9}{
            \makecell[{{p{\linewidth}}}]{
                \texttt{\tiny{[GM$|$GM]}}
                \texttt{game\_result = LOSE} \\
            }
        }
    }
    & & \\ \\

\end{supertabular}
}

\end{document}
